\section*{Chapter \ref{ch:b1:intermolecular-forces-solvents} --- Intermolecular Forces and Solvents}

\begin{solution}{exo:b1:intermolecular-forces-solvents:1}
Argon: London only. Hydrogen chloride: Keesom, Debye and London, London
dominating (\cref{exo:b1:intermolecular-forces-solvents:8}); chlorine is
too large for strong \omterm{def:b1:intermolecular-forces-solvents:hbond}{hydrogen bonds}. Water: \omterm{def:b1:intermolecular-forces-solvents:hbond}{hydrogen bonds} dominate, with
the three van der Waals terms. Methane: London only. Diiodine: London
only, strong because the molecule is large and very polarisable.
\end{solution}

\begin{solution}{exo:b1:intermolecular-forces-solvents:2}
Between their own molecules: \ce{CH3OH} (\ce{O-H}), \ce{CH3NH2}
(\ce{N-H}) and \ce{HF}. \ce{CH3OCH3} and \ce{CH3F} have \omterm{def:b1:lewis-resonance-vsepr:lewis}{lone pairs} on O
or F but no hydrogen on them: they can only accept \omterm{def:b1:intermolecular-forces-solvents:hbond}{hydrogen bonds}, from
water for instance. \ce{CH4} neither gives nor accepts.
\end{solution}

\begin{solution}{exo:b1:intermolecular-forces-solvents:3}
Water and ethanol: polar, protic. Propanone and acetonitrile: polar,
aprotic. Dichloromethane: weakly polar, aprotic. Hexane: apolar, aprotic.
\end{solution}

\begin{solution}{exo:b1:intermolecular-forces-solvents:4}
Diiodine is apolar: in hexane it exchanges \omterm{def:b1:intermolecular-forces-solvents:vdw}{London interactions} for \omterm{def:b1:intermolecular-forces-solvents:vdw}{London
interactions}, while in water it would break \omterm{def:b1:intermolecular-forces-solvents:hbond}{hydrogen bonds} without
replacing them. Sodium chloride is ionic: only a solvent of high
permittivity can separate its ions and only a polar one solvate them;
hexane ($\varepsilon_r = 1.89$) does neither.
\end{solution}

\begin{solution}{exo:b1:intermolecular-forces-solvents:5}
Methoxymethane has no \ce{O-H}: it cannot donate \omterm{def:b1:intermolecular-forces-solvents:hbond}{hydrogen bonds} and boils
lowest. Methanol and ethanol are both hydrogen-bonded; ethanol, with one
more \ce{CH2}, has stronger \omterm{def:b1:intermolecular-forces-solvents:vdw}{London interactions}. Ethanoic acid forms
hydrogen-bonded dimers and is the heaviest: it boils highest.
\end{solution}

\begin{solution}{exo:b1:intermolecular-forces-solvents:6}
In vacuum $F = e^2/(4\pi\varepsilon_0 r^2) = (1.602 \times
10^{-19})^2/(4\pi \times 8.854 \times 10^{-12} \times (5.00 \times
10^{-10})^2) = \qty{9.23e-10}{N}$. In water $F/78.4 = \qty{1.18e-11}{N}$;
in hexane $F/1.89 = \qty{4.88e-10}{N}$. Water weakens the attraction 41
times more than hexane: it can keep the ions apart, hexane cannot.
\end{solution}

\begin{solution}{exo:b1:intermolecular-forces-solvents:7}
Mixing propanone with water breaks some water--water \omterm{def:b1:intermolecular-forces-solvents:hbond}{hydrogen bonds} but
forms new ones between water and the oxygen of \ce{C=O}, plus
dipole--dipole interactions: the balance is favourable. Hexane offers
water only weak \omterm{def:b1:intermolecular-forces-solvents:vdw}{London interactions} in exchange for the \omterm{def:b1:intermolecular-forces-solvents:hbond}{hydrogen bonds} it
would break: the two liquids stay apart.
\end{solution}

\begin{solution}{exo:b1:intermolecular-forces-solvents:8}
\ce{HBr} has more electrons and a larger, more polarisable cloud: its
\omterm{def:b1:intermolecular-forces-solvents:vdw}{London interaction} exceeds that of \ce{HCl} by more than its \omterm{def:b1:intermolecular-forces-solvents:vdw}{Keesom
interaction} falls short. London dominates.
\end{solution}

\begin{solution}{exo:b1:intermolecular-forces-solvents:9}
\chemfig{CH_3-C(=[:60]O)-[:-60]O-[:0]H} facing its partner, two
$\ce{O-H}\cdots\ce{O=C}$ bonds closing an eight-membered ring. In benzene,
apolar and aprotic, nothing competes with these bonds: the acid is
dimerised and behaves as a particle of about \qty{120}{g/mol}. In water,
each acid molecule hydrogen-bonds with water instead: the dimers break up.
\end{solution}

\begin{solution}{exo:b1:intermolecular-forces-solvents:10}
Sodium chloride is made of ions: water, of high permittivity, separates
them (dissociation) and hydrates them (oxygen towards \ce{Na+}, hydrogen
towards \ce{Cl-}). Hydrogen chloride is a \omterm{def:b1:lewis-resonance-vsepr:dipole}{polar molecule}: water, polar,
turns the \ce{H-Cl} bond into an ion pair \ce{H3O+ Cl-} (ionisation),
separates the pair (dissociation) and hydrates the ions. Hexane is apolar
and of low permittivity: it neither ionises nor dissociates, so no ions
carry the current.
\end{solution}

\begin{solution}{exo:b1:intermolecular-forces-solvents:11}
$(174.1 - 36.1)/5 = \qty{27.6}{\celsius}$ per \ce{CH2}. The last
increment (nonane to decane) is \qty{23.4}{\celsius}, so undecane should
boil near $174 + 22 \approx \qty{196}{\celsius}$. Each \ce{CH2} adds the
same London contribution, but a smaller and smaller fraction of the
molecule's total, while the temperature needed to overcome the
interactions is already high.
% ledger: bp:C9H20
\end{solution}

\begin{solution}{exo:b1:intermolecular-forces-solvents:12}
\omterm{def:b1:intermolecular-forces-solvents:amphiphile}{Hydrophilic} part: the sulfate head \ce{-OSO3^-} with its \ce{Na+} counter-ion;
\omterm{def:b1:intermolecular-forces-solvents:amphiphile}{hydrophobic} part: the twelve-carbon chain. At the surface, heads in the
water and tails in the air; in a micelle, tails inside and heads outside.
The tails dissolve in the grease of the stain, the heads keep it in
contact with water, and the grease is lifted off in micelles.
\end{solution}

\begin{solution}{pb:b1:intermolecular-forces-solvents:1}
\textbf{1.} The \omterm{def:b1:intermolecular-forces-solvents:vdw}{London interaction}, the only one between apolar atoms.
\textbf{2.} From helium to xenon the number of electrons and the size grow,
hence the \omterm{def:b1:periodicity:polarisability}{polarisability} and the London attraction.
\textbf{3.} Helium has two electrons held tightly by its nucleus: the
least polarisable atom, with the weakest London attraction.
\textbf{4.} $68.8 - 36.1 = \qty{32.7}{\celsius}$: one \ce{CH2} adds
electrons and contact surface, hence London attraction.
\textbf{5.} London, growing with \omterm{def:b1:periodicity:polarisability}{polarisability}: $\ce{CH4} < \ce{SiH4} <
\ce{GeH4}$.
\textbf{6.} Carbon is not electronegative enough to polarise \ce{C-H}
strongly, and methane has no \omterm{def:b1:lewis-resonance-vsepr:lewis}{lone pair} to accept a \omterm{def:b1:intermolecular-forces-solvents:hbond}{hydrogen bond}.
\textbf{7.} Water: polar, protic, of the highest permittivity; it
separates and solvates both ions.
\textbf{8.} $78.4/1.89 = 41$: the attraction is 41 times weaker in water.
\textbf{9.} Acetonitrile (or propanone): polar enough to dissolve the salt,
aprotic so that the anion stays free.
\textbf{10.} Hexane: immiscible with water, apolar like diiodine. The iodine
passes into the hexane, the upper layer (hexane is less dense than water).
\textbf{11.} Ethanol is \omterm{def:b1:intermolecular-forces-solvents:miscibility}{miscible} with water: no second layer forms.
\textbf{12.} Its oxygen accepts \omterm{def:b1:intermolecular-forces-solvents:hbond}{hydrogen bonds} from water, but it cannot
donate any and its two ethyl groups are \omterm{def:b1:intermolecular-forces-solvents:amphiphile}{hydrophobic}: it dissolves only
slightly.
\textbf{13.} $-66.4 - (-85.0) = \qty{18.6}{\celsius}$ per \omterm{def:b1:periodicity:table}{period};
extrapolated \ce{HF}: $-85.0 - 18.6 = \qty{-103.6}{\celsius}$.
\textbf{14.} \ce{HF} boils at \qty{19.6}{\celsius}, \qty{123}{\celsius}
above the extrapolation.
\textbf{15.} $-62.5 - (-87.8) = \qty{25.3}{\celsius}$; extrapolated
\ce{NH3}: \qty{-113.1}{\celsius}; actual $-33.4$, \qty{80}{\celsius}
higher.
\textbf{16.} $-88.1 - (-112.0) = \qty{23.9}{\celsius}$; extrapolated
\ce{CH4}: \qty{-135.9}{\celsius}; actual $-161.5$, below. Methane forms no
\omterm{def:b1:intermolecular-forces-solvents:hbond}{hydrogen bonds}; it is simply small and weakly polarisable.
\textbf{17.} Water (about \qty{179}{\celsius}) $>$ \ce{HF} (123) $>$
\ce{NH3} (80).
\textbf{18.} \ce{NH3} can give three but accept only one (one \omterm{def:b1:lewis-resonance-vsepr:lewis}{lone pair});
\ce{HF} can accept three but give only one. Water gives two and accepts
two: every molecule can take part in four \omterm{def:b1:intermolecular-forces-solvents:hbond}{hydrogen bonds}, a full
three-dimensional network.
\textbf{19.} \ce{H2S} and \ce{H2Se} form no significant \omterm{def:b1:intermolecular-forces-solvents:hbond}{hydrogen bonds}:
their boiling points follow the London trend that water would follow
without them. Water itself is the anomaly being measured.
\textbf{20.} $-41.3 - (-60.3) = \qty{19.0}{\celsius}$ per \omterm{def:b1:periodicity:table}{period}.
\textbf{21.} $T(p) = -60.3 + 19.0\,(p - 3)$ in \unit{\celsius}.
\textbf{22.} $T(2) = -60.3 - 19.0 = \qty{-79.3}{\celsius}$.
\textbf{23.} $100.0 - (-79.3) = \qty{179}{\celsius}$.
\textbf{24.} Without its \omterm{def:b1:intermolecular-forces-solvents:hbond}{hydrogen bonds} water would boil near
\textbf{\qty{-79}{\celsius}}, and would be a gas on Earth.
\end{solution}
